%!TEX root = ../main.tex
% =============================================================
%  CAPÍTULO 2 — METODOLOGÍA             (Entrega 2 · 5 % · 5/5)
% =============================================================

\chapter{Metodología y procedimientos}\label{cap:metodologia}

\begin{guia}[Guía de la cátedra: qué se evalúa en este capítulo]
La cátedra exige \textbf{dos metodologías}. El error más común es
entregar solo una.
\begin{enumerate}
    \item \textbf{Metodología de la investigación}: área temática,
          planteamiento y delimitación del problema, marco teórico
          previsto, diseño concreto (fases), indicadores medibles con
          fuente y umbral, técnicas de recolección, instrumentos, datos,
          procesamiento, análisis crítico y síntesis. Tipo de
          investigación: en general, aplicada.
    \item \textbf{Metodología de dirección del proyecto}: marco de trabajo
          justificado, alcance y entregables, cronograma con fechas reales
          (Gantt con hitos que coinciden con las entregas), costos, riesgos
          (referencia al capítulo~\ref{cap:riesgos}), calidad, recursos y
          satisfacción de las partes interesadas.
    \item \textbf{Declaración de uso de IA}: qué tareas fueron asistidas y
          cómo se verificaron. Los prompts se guardan en una carpeta
          compartida con la cátedra.
\end{enumerate}
Extensión orientativa: 6 a 10 páginas. Referencias habituales:
\citet{hernandezsampieri2014} para investigación y \citet{pmi2021}
para gestión.
\end{guia}

\section{Área temática}

Este proyecto se desarrolla en el ámbito de soluciones de software
inteligentes aplicadas a la gestión integral de los eventos, con el
objetivo de entender las dificultades que existen al momento de
planificar y coordinar un evento, para así analizar de qué manera las
tecnologías de hoy en día pueden simplificar el proceso.

Por este motivo, el proyecto integra conceptos relacionados con
marketplaces digitales, sistemas de gestión comercial y operativa,
sistemas de recomendación con el uso de la \gls{ia} y Gemelos Digitales
para la simulación de escenarios.

\section{Planteamiento del problema}

Para organizar un evento social o empresarial hay que coordinar muchos
servicios, como proveedores del lugar, catering, fotografía, decoración,
ambiente, entretenimiento, proveedores de sonido e iluminación, además de
algunos adicionales que se requieren.

Hasta ahora, este tipo de servicios generalmente se contrataban por
separado a través de redes sociales, plataformas de mensajería, sitios web
y recomendaciones de boca en boca. Esto provoca que la información se
disperse por internet, y que los usuarios necesiten pasar mucho tiempo
buscando qué proveedores existen o comparando opciones y disponibilidad,
además de tomar decisiones en situaciones en las que disponen de datos
limitados.

Sin embargo, muchos usuarios contratan a un organizador de eventos que
actúa como nexo entre los distintos proveedores, agrega información y
trabaja con ellos en la planificación del evento. Aunque este enfoque
facilita la organización, conlleva gastos adicionales y obliga a depender
de plataformas de terceros.

Del mismo modo, los recintos y proveedores gestionan sus reservas,
presupuestos, pagos y disponibilidad mediante sistemas independientes o
procesos manuales, dejando pocas posibilidades de que todos los
participantes de un evento interactúen entre sí.

Esto lleva a explorar si un marketplace de servicios, junto con
herramientas de gestión y herramientas inteligentes de emparejamiento y
simulación, mejora la organización de eventos en general o ayuda en la
toma de decisiones.

\section{Delimitación de la investigación}

Este proyecto está enfocado en el desarrollo de una plataforma web y móvil
orientada a la organización de eventos sociales y empresariales.

La investigación está enfocada a personas que ya hayan organizado algún
tipo de evento, a los proveedores de productos y servicios relacionados
con la industria y a los propietarios de los salones, ya que son ellos
los que presentan las dificultades de llevar a cabo la organización y
planificación del evento.

La propuesta está delimitada a un estudio y desarrollo de un marketplace
de servicios, herramientas de gestión comercial y operativas, y
funcionalidades inteligentes de simulación con Gemelos Digitales y
recomendación de productos mediante \gls{ia}.

\section{Marco teórico previsto}

\begin{description}
    \item[Marketplaces digitales.] Se estudia el funcionamiento de los
          marketplaces digitales, desde sus funcionalidades hasta cómo
          permiten conectar al usuario con los proveedores, buscando
          encontrar las ventajas y limitaciones de las plataformas.
    \item[Sistemas de gestión comercial y operativa.] Se investigan los
          sistemas de gestión utilizados por los salones de fiestas,
          buscando cómo llevan a cabo los procesos de reserva, presupuesto,
          pagos, clientes y disponibilidad de recursos.
    \item[Sistemas de recomendación.] Se investigan los conceptos
          relacionados con la \gls{ia} y las tecnologías que permiten
          procesar información e identificar patrones que ayudan al
          usuario en la toma de decisiones, enfocándose en cómo estas
          herramientas pueden incorporarse dentro de una plataforma.
    \item[Inteligencia artificial.] Se estudian los conceptos relacionados
          con \gls{ia} y las tecnologías que permiten procesar información
          e identificar patrones para asistir en la toma de decisiones,
          analizando cómo pueden incorporarse funcionalidades inteligentes
          dentro de la plataforma propuesta.
    \item[\glspl{g:gemelo-digital}.] Se investiga el concepto de Gemelo
          Digital y sus aplicaciones actuales en diferentes industrias,
          con el objetivo de analizar la posibilidad de representar
          virtualmente un evento y simular diferentes escenarios antes de
          su realización.
    \item[Gestión integral de eventos.] Se estudian los procesos
          involucrados en la organización de eventos sociales y
          empresariales, incluyendo la contratación de proveedores, la
          coordinación de servicios, la administración de recursos y la
          planificación de actividades.
\end{description}

\section{Diseño de investigación}

Se define un tipo de investigación aplicada, desarrollada en cinco fases
sucesivas.

\subsection{Fase 1: Investigación y análisis de la industria}

Se realiza un relevamiento de la industria con el fin de comprender cómo
se organizan actualmente los eventos tanto sociales como empresariales, y
de qué manera las personas acceden a la contratación de productos y
servicios. Además, se analizan las diferentes plataformas existentes
relacionadas con la organización de eventos, tales como WeddingWire, The
Knot, Eventbrite, Cvent, Eventtia y Planning Pod, con el propósito de
identificar sus principales funcionalidades, limitaciones y oportunidades
de mejora.

\subsection{Fase 2: Relevamiento de requerimientos}

Para relevar los procesos actuales de organización y gestión de eventos,
identificar las principales dificultades y conocer las necesidades de
cada usuario, se realizan encuestas y entrevistas semiestructuradas
dirigidas a tres grupos principales: personas que hayan organizado
eventos, propietarios o responsables de salones y proveedores de
servicios. La información obtenida permite definir los requerimientos
funcionales y no funcionales del sistema.

\subsection{Fase 3: Definición y diseño de la solución}

Con los relevamientos se define cómo está compuesta la plataforma,
quiénes son los usuarios que participan en ella y de qué manera se
relacionan entre sí los clientes, los salones y los proveedores. Se
determina cómo funciona el marketplace y qué herramientas son necesarias
para administrar reservas, presupuestos y contrataciones, y se define la
incorporación de recomendaciones inteligentes y del Gemelo Digital.

\subsection{Fase 4: Desarrollo del prototipo funcional}

Se desarrolla una primera versión funcional de la plataforma web y móvil,
incorporando las principales funcionalidades definidas en las fases
anteriores, incluyendo el marketplace, las herramientas de gestión para
salones y proveedores, y las funcionalidades inteligentes de
recomendación y simulación.

\subsection{Fase 5: Validación funcional y análisis de resultados}

Se realizan pruebas con usuarios de los tres grupos involucrados en la
investigación, comparando los datos obtenidos con la información relevada
en la Fase 2, para determinar si la propuesta logra resolver o simplificar
la problemática identificada.

\section{Instrumentos de recolección de datos}

Para obtener la información necesaria se utilizan distintos instrumentos
según las características de cada grupo. Para las personas que hayan
organizado eventos, propietarios de salones y proveedores, se utilizan
formularios de encuestas desarrollados mediante herramientas digitales,
principalmente Google Forms. Además, se realizan comparaciones entre las
distintas plataformas y herramientas que existen actualmente en la
industria, permitiendo identificar oportunidades de mejora y aspectos
diferenciales (apéndice~\ref{apx:encuesta}).

\section{Procesamiento de los datos}

Una vez recolectada la información, los datos se organizan y separan
según el tipo de participante (clientes, salones y proveedores), y se
registran en planillas y tablas que permiten ordenar las respuestas,
identificar dificultades comunes y definir las funcionalidades más
importantes a incorporar en la plataforma.

\section{Análisis de los datos}

Los datos obtenidos se analizan mediante tablas, gráficos y estadísticas
descriptivas. La información de las entrevistas se analiza de forma
cualitativa, con el propósito de comprender las problemáticas actuales,
identificar necesidades de gestión y detectar oportunidades de mejora
para el diseño de la plataforma (apéndice~\ref{apx:entrevista}).

\section{Síntesis y conclusiones}

Se realiza una síntesis de la información obtenida durante la
investigación, analizando el nivel de interés de clientes, salones y
proveedores en el proyecto, y definiendo las funcionalidades que debe
incorporar la plataforma para responder de manera más eficiente a las
necesidades actuales de la industria de eventos.

\section{Declaración de uso de \acrshort{ia}}

\completar{Qué tareas del capítulo fueron asistidas por herramientas de
\gls{ia} y cómo se verificó cada resultado (ver
declaración completa en el preliminar correspondiente).}

\section{Metodología de dirección del proyecto}

\subsection{Alcance}

El proyecto incluye el análisis, diseño, desarrollo y validación
funcional de una plataforma web y móvil orientada a la organización
integral de eventos sociales y empresariales. La solución contempla un
marketplace para buscar, comparar y contratar servicios (salones,
catering, fotografía, decoración, ambientación, entretenimiento y otros
proveedores), herramientas de gestión para salones y proveedores, y
funcionalidades inteligentes de recomendación y Gemelo Digital para
simular distintos escenarios antes de la realización del evento.

\subsection{Cronograma}

Los tiempos indicados representan una estimación inicial y se ajustan
según el avance de la investigación y el desarrollo del proyecto. Se
estima una duración total de aproximadamente diez meses, organizada en
cinco etapas:

\begin{itemize}
    \item \textbf{Etapa 1 — Investigación y relevamiento (junio)}: estudio
          del estado del arte, investigación de plataformas existentes y
          relevamiento mediante encuestas y entrevistas a clientes,
          salones y proveedores.
    \item \textbf{Etapa 2 — Análisis y diseño de la solución
          (julio-agosto)}: definición de requerimientos funcionales y no
          funcionales, arquitectura general del sistema, modelo de datos
          y diseño conceptual del marketplace, las herramientas de
          gestión y el Gemelo Digital.
    \item \textbf{Etapa 3 — Desarrollo de la plataforma
          (agosto-septiembre)}: desarrollo del marketplace de servicios,
          las herramientas de administración y las funcionalidades
          inteligentes de recomendación y simulación.
    \item \textbf{Etapa 4 — Validación funcional y pruebas
          (octubre-noviembre)}: pruebas funcionales y evaluaciones con
          usuarios representativos, identificando mejoras y ajustes
          necesarios.
    \item \textbf{Etapa 5 — Documentación y cierre (diciembre)}:
          elaboración de la documentación final, análisis de resultados y
          preparación de la presentación final.
\end{itemize}

Este cronograma de etapas de investigación se complementa con el
calendario oficial de entregas de la cátedra (figura~\ref{fig:gantt}),
que fija las fechas concretas de cada capítulo.

\begin{figure}[H]
    \centering
    \begin{ganttchart}[
        hgrid, vgrid={*{6}{draw=none}, dotted},
        x unit=0.42mm,
        y unit chart=6mm, y unit title=7mm,
        time slot format=isodate,
        bar label font=\small, milestone label font=\small,
        title label font=\small,
        bar/.append style={fill=mainColor!60, draw=none},
        milestone/.append style={fill=mainColor, draw=none},
    ]{2026-03-16}{2026-11-20}
        \gantttitlecalendar{month=shortname} \\
        \ganttbar{Idea y resumen}{2026-03-16}{2026-04-10} \\
        \ganttbar{Metodología}{2026-04-06}{2026-05-05} \\
        \ganttbar{Introducción}{2026-05-06}{2026-06-09} \\
        \ganttbar{Marco teórico}{2026-06-01}{2026-07-28} \\
        \ganttbar{Marco tecnológico}{2026-07-01}{2026-08-25} \\
        \ganttbar{Justif. económica}{2026-08-10}{2026-09-15} \\
        \ganttbar{Riesgos}{2026-09-01}{2026-09-29} \\
        \ganttbar{Desarrollo del sistema}{2026-06-15}{2026-10-10} \\
        \ganttbar{Resultados}{2026-09-20}{2026-10-13} \\
        \ganttbar{Conclusiones}{2026-10-05}{2026-10-20} \\
        \ganttbar{Integración y revisión}{2026-10-21}{2026-11-16} \\
        \ganttmilestone{Entrega final}{2026-11-17}
    \end{ganttchart}
    \caption{Cronograma del proyecto (ciclo lectivo 2026)}
    \label{fig:gantt}
\end{figure}

\subsection{Costos del proyecto}

Esta sección corresponde a los costos de llevar adelante el proyecto de
tesis en sí --- no a la factibilidad económica del negocio una vez
lanzada la plataforma---: el tiempo, las herramientas y los eventuales
recursos necesarios para investigar, diseñar y desarrollar el prototipo
dentro del plazo académico.

\begin{table}[H]
    \centering
    \caption{Costos del proyecto de tesis (estimado)}
    \label{tab:costos-tesis}
    \small
    \begin{tabularx}{\textwidth}{L L L}
        \toprule
        \tb{Concepto} & \tb{Descripción} & \tb{Costo estimado} \\
        \midrule
        Tiempo del alumno & Dedicación estimada del tesista durante los 10 meses del proyecto (no representa un desembolso, pero es el principal recurso invertido) & Sin costo monetario directo \\
        Herramientas de desarrollo & Editor de código, control de versiones (Git/GitHub en su plan gratuito), entornos de desarrollo local & USD 0 (herramientas gratuitas) \\
        Encuestas y formularios & Google Forms para las encuestas de la Fase 2 & USD 0 (plan gratuito) \\
        Hosting de prueba / dominio de desarrollo & Entorno de staging para mostrar el prototipo funcional durante la defensa & USD 5 -- 15 / mes \\
        Incentivos a entrevistados (opcional) & Eventual compensación simbólica a salones/proveedores que dediquen tiempo a entrevistas extensas & A definir según disponibilidad \\
        Impresión y encuadernación & Copias físicas del documento final, si la cátedra las requiere & USD 20 -- 40 (estimado) \\
        \bottomrule
    \end{tabularx}
\end{table}

\subsection{Recursos}

Además de los recursos preliminares descriptos en el
capítulo~\ref{cap:introduccion}, la metodología de gestión del proyecto
requiere detallar los recursos concretos que se utilizan durante las
distintas etapas:

\begin{itemize}
    \item \textbf{Recursos humanos}: el alumno, como responsable del
          análisis, diseño, desarrollo y validación; el tutor/profesor de
          la cátedra, para el seguimiento y revisión periódica del
          avance; y los participantes de las entrevistas y encuestas
          (clientes, salones y proveedores, incluyendo Rodríguez Catering
          como caso de estudio real).
    \item \textbf{Recursos materiales y tecnológicos}: una computadora de
          desarrollo, conexión a internet, las herramientas de software
          detalladas en el capítulo~\ref{cap:marco-tecnologico} (React,
          Laravel, MySQL/SQLite), un repositorio Git para el control de
          versiones, y Google Forms para las encuestas de la Fase 2.
    \item \textbf{Recursos de tiempo}: detallados en el diagrama de Gantt
          de la figura~\ref{fig:gantt}.
\end{itemize}

\subsection{Riesgos}

Se identifican los siguientes riesgos principales del proyecto de tesis
en sí --- distintos de los riesgos del negocio--- y sus estrategias de
mitigación:

\begin{itemize}
    \item \textbf{Riesgo 1 — Baja participación en las encuestas y
          entrevistas}: la obtención de información proveniente de
          clientes, salones y proveedores resulta fundamental para
          comprender las necesidades reales del sector.
          \emph{Mitigación}: realizar el relevamiento con distintos
          perfiles de participantes y utilizar diferentes medios de
          contacto.
    \item \textbf{Riesgo 2 — Dificultad para representar adecuadamente el
          comportamiento de los eventos}: la implementación de
          herramientas de recomendación y simulación requiere comprender
          cómo se desarrollan los eventos. \emph{Mitigación}: trabajar con
          información de las entrevistas y validar los escenarios con
          usuarios vinculados a la industria.
    \item \textbf{Riesgo 3 — Complejidad de integrar diferentes
          funcionalidades en una misma plataforma}: la solución reúne un
          marketplace, herramientas de gestión y funcionalidades
          inteligentes que deben funcionar de manera integrada.
          \emph{Mitigación}: diseñar una arquitectura modular que permita
          desarrollar y validar cada componente de manera independiente.
    \item \textbf{Riesgo 4 — Cambios en los requerimientos identificados
          durante la investigación}: podrían surgir nuevas necesidades a
          medida que avance el relevamiento. \emph{Mitigación}: mantener
          un proceso de revisión continua y priorizar las funcionalidades
          más importantes.
    \item \textbf{Riesgo 5 — Limitaciones de tiempo durante el
          desarrollo}: la cantidad de funcionalidades propuestas podría
          generar dificultades para completar todas las etapas.
          \emph{Mitigación}: definir prioridades de desarrollo y
          establecer entregas parciales.
\end{itemize}

\subsection{Calidad}

Se implementan mecanismos de verificación y validación a lo largo de todo
el proyecto: pruebas sobre las principales características de la
plataforma (marketplace, herramientas de gestión, recomendación y
simulación), pruebas de integración entre componentes, pruebas de
usabilidad con los tres grupos de usuarios, verificación del cumplimiento
de los requerimientos funcionales y no funcionales, y elaboración de la
documentación técnica y funcional del sistema.

\subsection{Satisfacción}

La satisfacción de los usuarios finales se evalúa mediante encuestas y
pruebas realizadas con personas que hayan organizado eventos, propietarios
de salones y proveedores de servicios, considerando: facilidad de uso y
claridad de la interfaz; utilidad del marketplace, de las herramientas de
gestión, de las recomendaciones y del Gemelo Digital; percepción de
mejora en la organización y toma de decisiones; satisfacción general;
intención de uso futuro; y disposición a recomendar la plataforma.
